%% RESUMO NA LINGUA VERNACULA (elemento obrigatorio -- NBR 6028)
%%
%% Regras:
%%   - paragrafo unico, sem recuo, justificado;
%%   - Times New Roman 12, espacamento simples;
%%   - entre 150 e 500 palavras;
%%   - verbo na voz ativa e na terceira pessoa do singular;
%%   - as palavras-chave sao definidas em config/dados.tex.
%%
%% Um bom resumo informa: objetivo geral; principais conceitos e autores da
%% fundamentacao teorica; etapas e acoes da execucao; populacao envolvida e
%% periodo; principais resultados alcancados; conclusoes.

\begin{resumo}
Esta dissertação investiga o uso de inteligência artificial no monitoramento
contínuo e na conformidade com a Lei Geral de Proteção de Dados (LGPD) no
Sistema Eletrônico de Informações (SEI). Parte do grande volume de processos
que tramitam no sistema e das limitações do controle manual, pouco escalável e
sujeito a falhas. O problema de pesquisa é como soluções de inteligência
artificial e de processamento de linguagem natural podem automatizar a
identificação de dados pessoais e sensíveis nos documentos. Os critérios para o
uso dessas técnicas são validados por especialistas pelo Método DELPHI. A pesquisa busca,
com isso, contribuir para a melhoria da segurança e da eficiência da gestão
documental pública, por meio de um fluxo que vai da leitura dos documentos do
SEI à emissão de alertas de conformidade.
\end{resumo}
